% Condensed from the source report, Sections 1 and 7.
\section{Introduction}
\label{sec:introduction}
Across meetings, phone calls, and continuously listening applications, a recognizer may need to transcribe one person despite changes in acoustics and speaker activity.
Target-speaker automatic speech recognition (TS-ASR) receives a short enrollment recording of that person and a recording to transcribe, called the mixture.
It must recognize the target's words despite competing voices, but also return an empty transcript when only other people speak.
A system that recognizes speech accurately yet frequently transcribes the wrong person fails this task.

Common evaluations such as WSJ0-2mix, LibriMix, and LibriSpeechMix overlay read utterances and assume that the target is present~\cite{hershey2016deep,cosentino2020librimix,kanda2020serialized}.
They measure recognition or separation under overlap, but cannot expose false alarms when the enrolled speaker is absent.
Prior target speech extraction work has studied this failure explicitly~\cite{delcroix2022listen,borsdorf2021universal,zhang2023speaker}.
Enrollment-conditioned speech models, including SQ-Whisper and CocktailASR, extend recognition to competing speakers~\cite{guo2024sqwhisper,zhang2025thinking,zhang2026xiaomi}.
Real meeting benchmarks address speaker-attributed transcription~\cite{liang2023second,vinnikov2024notsofar}, but this differs from recognizing a single enrolled person and rejecting all others.
The remaining evaluation gap is to test both behaviors across conversational timing, degraded enrollment, and real recordings.

We introduce \model (One For All Target-Speaker ASR), using the same model parameters across noise, reverberation, turn-taking, overlap, short utterances, and target absence.
It builds on Qwen3-ASR-1.7B~\cite{shi2026qwen3asr}, conditions its audio encoder on enrollment speech, and trains with target-absent counterparts across scene types.
This extends the speaker-selection capability studied in AmphionASR~1.0~\cite{amphion2026amphionasr}.
To evaluate this scenario coverage, we introduce \bench, with 24{,}101 Mandarin samples in 248 subsets.
It combines controlled synthetic conditions with real meetings, telephone calls, and web conversations, measuring recognition, complete misses, and false alarms together.

The benchmark changes how systems compare.
When only the target speaks in real recordings, \model and Xiaomi-CocktailASR-1 have nearly identical character error rates (CERs), 7.5\% and 7.3\%.
When others speak while the target is absent, their false-alarm rates are 2.6\% and 96.1\%.
This contrast motivates evaluating target absence as a distinct requirement rather than treating silence-only tests as sufficient.
